% !TEX root = ../../../main.tex
% Sources: LEC02, IMG_9118--IMG_9119.
\section{Convergence and iteration}
\label{sec:analysis-i:convergence-preview}

\newthought{A sequence may approach a number that its terms never reach.}
This is why a field containing all the rational numbers can still be
missing limits. We now make the meaning of ``approach'' precise.

\begin{definition}
  A sequence $(x_n)$ of real numbers \emph{converges} to $L\in\R$ if,
  for every $\varepsilon>0$, there exists $N\in\N$ such that
  \[
    n\geq N\quad\Longrightarrow\quad |x_n-L|<\varepsilon.
  \]
  We write $x_n\to L$ or $\lim_{n\to\infty}x_n=L$.
\end{definition}
The index $N$ may depend on $\varepsilon$. After that index, every term
lies in the $\varepsilon$-neighborhood $(L-\varepsilon,L+\varepsilon)$.
For convergence \emph{in $\Q$}, the limit must belong to $\Q$ as well.
The rational partial sums defining $e$ converge in $\R$, but not in $\Q$.
\marginnote{The absolute value $|x|$ is $x$ for $x\geq0$ and $-x$ otherwise.
It measures distance from zero. Its elementary inequalities are proved in
\S~\ref{sec:analysis-i:elementary-functions}.}

We will use two elementary consequences of this definition. Limits are
unique, since if $L\ne M$, the disjoint neighborhoods of radius
$|L-M|/3$ cannot both contain all sufficiently late terms. Also, if
$x_n\leq y_n$, $x_n\to x$, and $y_n\to y$, then $x\leq y$.
Otherwise neighborhoods of radius $(x-y)/3$ contradict $x_n\leq y_n$.

\subsection{Newton's idea}
Suppose we want to solve $f(x)=0$. At a current approximation $x_n$ with
$f'(x_n)\ne0$, replace the graph by its tangent line
\[
  y=f(x_n)+f'(x_n)(x-x_n).
\]
Its horizontal intercept supplies the next approximation,
\begin{equation}
  x_{n+1}=x_n-\frac{f(x_n)}{f'(x_n)}.
  \label{eq:analysis-i:newton}
\end{equation}
Figure~\ref{fig:newton-square-root} illustrates one step. This is a
motivation from calculus; a general convergence theorem for Newton's
method is not being assumed.
\begin{marginfigure}
  \centering
  \begin{tikzpicture}[analysis diagram,x=14mm,y=7mm]
  \draw[->] (-0.15,0) -- (2.8,0) node[right] {$x$};
  \draw[->] (0,-2.5) -- (0,4.2) node[above] {$y$};
  \draw[analysis curve,->,domain=0:2.4,samples=60]
    plot (\x,{\x*\x-2});
  \draw[gray,domain=1:2.5] plot (\x,{4*\x-6});
  \draw[analysis guide] (2,0) -- (2,2);
  \fill (2,2) circle[radius=1.6pt];
  \fill (1.414214,0) circle[radius=1.6pt];
  \fill (1.5,0) circle[radius=1.6pt];
  \fill (0,-2) circle[radius=1.6pt];
  \node[left] at (0,-2) {$-2$};
  \node[below left,analysis label] at (0,0) {$0$};
  \node[above left,analysis label] at (1.414214,0) {$\sqrt2$};
  \node[below,analysis label] at (1.7,-0.2) {$\tfrac32$};
  \draw[gray] (1.5,-0.1) -- (1.7,-0.35);
  \node[below,analysis label] at (2,0) {$2$};
  \node[left,analysis label] at (1.9,2.9) {$x^2-2$};
\end{tikzpicture}

  \caption[One Newton step for a square root.]{For $f(x)=x^2-2$, the tangent at $x_0=2$ meets the axis at
    $x_1=3/2$. The parabola is shown for $x\geq0$.}
  \label{fig:newton-square-root}
\end{marginfigure}

For $f(x)=x^2-a$, where $a>0$, equation~\eqref{eq:analysis-i:newton} becomes
\[
  x_{n+1}=\frac12\left(x_n+\frac{a}{x_n}\right).
\]
A positive limit $L$, if it exists, must satisfy $L^2=a$. A fixed-point
calculation alone does not establish that such a limit exists.

Here is a convergence argument using two results proved later in this
chapter. Let $r=\sqrt a>0$, whose existence follows from
Theorem~\ref{thm:analysis-i:positive-roots}, and start with $x_0>0$.
Every iterate is positive, and
\[
  x_{n+1}-r=\frac{(x_n-r)^2}{2x_n}\geq0.
\]
Thus $x_n\geq r$ for $n\geq1$, and then
$x_{n+1}-x_n=(a-x_n^2)/(2x_n)\leq0$.
Monotone convergence gives a limit $L\geq r>0$. Passing to the limit in
$2x_nx_{n+1}=x_n^2+a$ gives $L^2=a$, hence $L=r$ by uniqueness of the
positive root. If $a$ and $x_0$ are rational, every iterate is rational;
for prime $a$, the limit is irrational.

\subsection{What kind of limit can a recurrence produce?}
The lecture also considered
\[
  x_{n+1}=a+\frac{1}{x_n},\qquad a\in\N,\quad x_0>0.
\]
If this sequence converges, its limit is positive because $x_n>a$ for
$n\geq1$. The limit therefore satisfies
\[
  L^2-aL-1=0,\qquad L=\frac{a+\sqrt{a^2+4}}{2}.
\]
It is algebraic and irrational. Indeed, the rational root theorem would
force a rational root of this monic polynomial to be an integer dividing
$1$, and neither $1$ nor $-1$ is a root for $a\geq1$.
Repeated substitution also suggests a continued fraction. This example
identifies a possible limit; it does not by itself prove convergence.

Limits of rational sequences need not be algebraic. The factorial partial
sums approach $e$. Another lecture example is Euler's identity
\[
  \sum_{n=1}^{\infty}\frac{1}{n^2}=\frac{\pi^2}{6},
\]
stated here without proof~\cite{Waldschmidt2017}.
The number $\pi^2/6$ is transcendental, since a nonzero rational polynomial
vanishing there would, after substituting $X^2/6$, give a nonzero rational
polynomial vanishing at $\pi$.

\subsection{Exercise on Newton's method}
We now apply Newton's method to the fixed-point equation for the continued
fraction above. As in the square-root example, the convergence argument
uses the monotone convergence theorem proved in
\S~\ref{sec:analysis-i:completeness-equivalences}.

% !TEX root = ../../hw02.tex
% Homework 02, Exercise 6. Shared problem statement.
\begingroup
\renewcommand{\labelenumi}{(\alph{enumi})}
\begin{exercise}
  \label{ex:hw02:newton-iteration}
  \solutionlink{hw02:newton-iteration}
  In class, we discussed the iteration process
  \[
    x_{n+1}=a+\frac{1}{x_n},\quad n=0,1,2,\ldots,
    \quad a\geq 1,\quad x_0>0,
  \]
  whose limit gives us the continuous fraction
  \[
    b=a+\frac{1}{a+\frac{1}{a+\frac{1}{a+\cdots}}}.
  \]
  That is, $b$ is a fixed point of the function
  \[
    f(x)=a+\frac{1}{x},\quad x>0,
  \]
  which is realized as a root of the function $F(x)=x-f(x)$ $(x>0)$.
  \begin{enumerate}
    \item Show that the iterative algorithm
      \[
        x_{n+1}=\frac{ax_n^2+2x_n}{1+x_n^2},
        \quad n=0,1,2,\ldots,\quad x_0>0,
      \]
      arises from Newton's method that may be used to approach $b$ of our problem.
    \item Assume that $\{x_n\}$ in (a) converges.
      Verify that the limit is indeed the continuous fraction $b$ defined above.
    \item Take $a=\frac{3}{2}$. Prove $0<x_n\leq 2$ for $n=1,2,\ldots$.
    \item Prove that $\{x_n\}_{n=1,2,\ldots}$ increases for any $x_0\ne 2$
      and is constant when $x_0=2$.
    \item Show that $\{x_n\}$ converges to $2$ with any initial state $x_0$
      (thereby establishing the ``global convergence'' of the algorithm).
  \end{enumerate}
\end{exercise}
\endgroup


\subsection{Exercise on Fibonacci ratios}
The next recurrence connects an integer sequence with bounded ratios.
As in the continued-fraction example, the last part asks for a possible
limit under the assumption that the ratios converge.

% !TEX root = ../../hw03.tex
% Source: HW3 Intro to Math Analysis I (1).pdf, page 1, question 4.
% "Increasing" in part (a) is non-strict, since F_1 = F_2 = 1.
\begingroup
\renewcommand{\labelenumi}{(\alph{enumi})}
\begin{exercise}
  \label{ex:hw03:fibonacci-ratios}
  \solutionlink{hw03:fibonacci-ratios}
  The Fibonacci sequence $\{F_n\}$ is defined by
  \[
    F_0=0,\quad F_1=1,\quad F_n=F_{n-2}+F_{n-1},
    \qquad n=2,3,\ldots.
  \]
  \begin{enumerate}
    \item Show that $\{F_n\}$ is an increasing sequence.
      \emph{Here increasing means nondecreasing.}
    \item Show that $\{F_n\}$ is unbounded.
    \item Show that the Fibonacci ratios
      \[
        r_n=\frac{F_n}{F_{n-1}},\qquad n=2,3,\ldots,
      \]
      is a bounded sequence.
    \item Assume the limit of $\{r_n\}$ exists. Find the limit
      (which happens to be the golden ratio).
  \end{enumerate}
\end{exercise}
\endgroup

